\documentclass[12pt,a4paper]{article}

\usepackage[T2A]{fontenc}
\usepackage[utf8]{inputenc}
\usepackage[russian]{babel}
\usepackage{amsmath,amssymb}
\usepackage{graphicx}
\usepackage{geometry}
\usepackage{listings}
\usepackage{xcolor}
\usepackage{hyperref}
\usepackage{caption}
\usepackage{float}

\geometry{left=2.5cm,right=2.5cm,top=2.5cm,bottom=2.5cm}

\lstset{
  basicstyle=\ttfamily\small,
  breaklines=true,
  frame=single,
  backgroundcolor=\color{gray!10},
  keywordstyle=\color{blue},
  commentstyle=\color{green!50!black},
  stringstyle=\color{red!70!black},
  language=Python,
  showstringspaces=false,
}

\title{Проект: Многослойный персептрон}
\author{Сагдуллин Амир 09 - 321}
\date{\today}

\begin{document}
\maketitle

\tableofcontents
\newpage

\section{Введение}

Цель проекта - реализовать многослойный персептрон (MLP) с нуля на Python,
без использования готовых библиотек машинного обучения (PyTorch, TensorFlow, Keras).
В работе применяются только NumPy для линейной алгебры, pandas для загрузки данных
и matplotlib для визуализации.

Задача - бинарная классификация объектов из датасета Breast Cancer Wisconsin
Diagnostic: по 30 числовым признакам, описывающим ядро клетки, необходимо определить,
является ли опухоль злокачественной (M) или доброкачественной (B).

\section{Первичный анализ данных (EDA)}

\subsection{Общая информация о датасете}

Датасет содержит 569 объектов и 32 столбца: идентификатор, метку класса
(диагноз) и 30 числовых признаков. Признаки описывают три группы характеристик
ядра клетки: средние значения (mean), стандартные ошибки (error) и худшие
значения (worst) для десяти базовых параметров.

\begin{figure}[H]
  \centering
  \includegraphics[width=\textwidth]{eda_head.png}
  \caption{Первые строки датасета}
\end{figure}

\subsection{Проверка качества данных}

\begin{figure}[H]
  \centering
  \includegraphics[width=\textwidth]{eda_info.png}
  \caption{Проверка пропусков, дубликатов и типов данных}
\end{figure}

Пропусков нет, полностью дублирующихся строк нет, дублирующихся
идентификаторов нет. Данные не требуют дополнительной очистки.

\subsection{Баланс классов}

\begin{figure}[H]
  \centering
  \includegraphics[width=0.7\textwidth]{eda_balance.png}
  \caption{Распределение классов}
\end{figure}

Класс B (доброкачественные) - 357 объектов (62.7\%), класс M (злокачественные) -
212 объектов (37.3\%). Наблюдается умеренный дисбаланс, поэтому при разделении
данных используется стратификация.

\subsection{Описательная статистика}

Признаки имеют существенно разные масштабы: например, \texttt{mean area}
варьируется в диапазоне 143.5-2501, а \texttt{mean smoothness} - в диапазоне
0.053-0.163. Это делает необходимой нормализацию перед обучением.

\subsection{Визуальный анализ признаков}

\begin{figure}[H]
  \centering
  \includegraphics[width=\textwidth]{eda_boxplot.png}
  \caption{Boxplot топ-5 признаков по корреляции с целевой переменной}
\end{figure}

Наиболее информативные признаки чётко разделяют классы: медиана класса M
значительно выше, чем класса B. Встречаются выбросы, но они отражают реальные
особенности данных (аномально крупные или мелкие ядра клеток), а не ошибки
измерений, поэтому удалять их нецелесообразно.

\begin{figure}[H]
  \centering
  \includegraphics[width=\textwidth]{eda_histograms.png}
  \caption{Распределения значений трёх наиболее информативных признаков по классам}
\end{figure}

Гистограммы подтверждают результаты корреляционного анализа: для каждого из трёх
признаков пики распределений классов B и M смещены относительно друг друга,
что говорит о высокой разделяющей способности этих характеристик. При этом
наблюдается перекрытие в средней части распределений - именно объекты из этой
зоны и становятся сложными случаями для классификатора.

\subsection{Корреляционный анализ}

\begin{figure}[H]
  \centering
  \includegraphics[width=\textwidth]{eda_corr_matrix.png}
  \caption{Корреляционная матрица признаков}
\end{figure}

\begin{figure}[H]
  \centering
  \includegraphics[width=0.8\textwidth]{eda_corr_target.png}
  \caption{Корреляция признаков с целевой переменной}
\end{figure}

Обнаружено 15 пар признаков с коэффициентом корреляции по модулю больше 0.95.
Основные тройки: \texttt{radius}, \texttt{perimeter}, \texttt{area} -
в каждой из групп (mean/error/worst) они почти линейно связаны между собой.
Для MLP мультиколлинеарность не критична, поэтому удаление признаков не выполнялось.

Наибольшая корреляция с целевой переменной: \texttt{worst concave points} (0.794),
\texttt{worst perimeter} (0.783), \texttt{mean concave points} (0.777).

\section{Разделение данных}

Данные разделены на обучающую и валидационную выборки в пропорции 80/20
со стратификацией по классу:

\begin{verbatim}
Train: shape=(456, 32)  B=286 (62.7%)  M=170 (37.3%)
Valid: shape=(113, 32)  B=71  (62.8%)  M=42  (37.2%)
\end{verbatim}

Стратификация позволила сохранить пропорцию классов в обеих выборках.
Случайное перемешивание выполнено с фиксированным seed = 42 для
воспроизводимости результатов.

\section{Архитектура нейронной сети}

\subsection{Прямой проход}

Для слоя с весами $W \in \mathbb{R}^{n_{\text{in}} \times n_{\text{out}}}$
и смещением $b \in \mathbb{R}^{1 \times n_{\text{out}}}$ прямой проход:

\begin{align}
z &= x W + b, \\
a &= f(z),
\end{align}

где $f$ — функция активации. В скрытых слоях используется сигмоида:

\begin{equation}
\sigma(z) = \frac{1}{1 + e^{-z}},
\end{equation}

на выходном слое — softmax для двух классов:

\begin{equation}
\text{softmax}(z)_i = \frac{e^{z_i}}{\sum_{j=1}^{2} e^{z_j}}.
\end{equation}

\subsection{Функция потерь}

Используется бинарная кросс-энтропия (для one-hot представления двух классов):

\begin{equation}
L = -\frac{1}{N} \sum_{i=1}^{N} \sum_{k=1}^{2} y_{ik} \log \hat{p}_{ik},
\end{equation}

где $y_{ik}$ — целевое значение, $\hat{p}_{ik}$ — предсказанная вероятность.

\subsection{Обратный проход}

Ключевое преимущество комбинации softmax + кросс-энтропия — упрощение градиента
на выходном слое:

\begin{equation}
\frac{\partial L}{\partial z^{(\text{out})}} = \frac{1}{N} (\hat{p} - y).
\end{equation}

Для скрытых слоёв с сигмоидой:

\begin{equation}
\frac{\partial L}{\partial z^{(l)}}
= \frac{\partial L}{\partial a^{(l)}} \odot a^{(l)} \odot (1 - a^{(l)}),
\end{equation}

\begin{equation}
\frac{\partial L}{\partial W^{(l)}} = (a^{(l-1)})^\top \frac{\partial L}{\partial z^{(l)}},
\quad
\frac{\partial L}{\partial b^{(l)}} = \sum \frac{\partial L}{\partial z^{(l)}}.
\end{equation}

\subsection{Инициализация весов}

Используется инициализация He:

\begin{equation}
W_{ij} \sim \mathcal{N}\left(0, \sqrt{\frac{2}{n_{\text{in}}}}\right).
\end{equation}

\section{Реализация}

Проект состоит из четырёх модулей:

\begin{itemize}
  \item \texttt{src/layers.py} - класс \texttt{DenseLayer} и функции активации
        (\texttt{sigmoid}, \texttt{relu}, \texttt{softmax});
  \item \texttt{src/model.py} - сборка сети (\texttt{createNetwork}), обучение
        (\texttt{fit}), предсказание (\texttt{predict}), сохранение и загрузка
        модели;
  \item \texttt{src/split\_data.py} - разделение данных;
  \item \texttt{src/train.py}, \texttt{src/predict.py} - командные интерфейсы
        для обучения и прогнозирования.
\end{itemize}

Архитектура сети задаётся через аргументы командной строки, что позволяет
легко экспериментировать с числом и размером слоёв.

\section{Обучение}

Гиперпараметры:

\begin{itemize}
  \item Архитектура: $30 \to 24 \to 24 \to 24 \to 2$
  \item Активации: sigmoid, sigmoid, sigmoid, softmax
  \item Функция потерь: бинарная кросс-энтропия
  \item Оптимизатор: SGD
  \item Learning rate: 0.0314
  \item Batch size: 8
  \item Количество эпох: 84
  \item Seed: 42
\end{itemize}

Команда запуска:

\begin{lstlisting}
python src/train.py --layer 24 24 24 --epochs 84 --loss binaryCrossentropy \
                    --batch_size 8 --learning_rate 0.0314
\end{lstlisting}

Лог обучения (фрагмент):

\begin{verbatim}
x_train shape : (456, 30)
x_valid shape : (113, 30)
epoch 01/84 - loss: 0.6558 - val_loss: 0.6572 - acc: 0.6272 - val_acc: 0.6283
epoch 10/84 - loss: 0.3080 - val_loss: 0.3244 - acc: 0.9364 - val_acc: 0.9381
epoch 20/84 - loss: 0.1054 - val_loss: 0.1209 - acc: 0.9781 - val_acc: 0.9735
epoch 30/84 - loss: 0.0710 - val_loss: 0.0925 - acc: 0.9868 - val_acc: 0.9735
epoch 40/84 - loss: 0.0588 - val_loss: 0.0872 - acc: 0.9868 - val_acc: 0.9735
epoch 50/84 - loss: 0.0525 - val_loss: 0.0859 - acc: 0.9890 - val_acc: 0.9735
epoch 60/84 - loss: 0.0483 - val_loss: 0.0867 - acc: 0.9912 - val_acc: 0.9646
epoch 70/84 - loss: 0.0453 - val_loss: 0.0909 - acc: 0.9912 - val_acc: 0.9646
epoch 80/84 - loss: 0.0429 - val_loss: 0.0911 - acc: 0.9912 - val_acc: 0.9558
epoch 84/84 - loss: 0.0421 - val_loss: 0.0943 - acc: 0.9912 - val_acc: 0.9558
> saving model 'saved_model.npy' to disk...
> learning curves saved to 'results/figures/learning_curves.png'
\end{verbatim}

\begin{figure}[H]
  \centering
  \includegraphics[width=\textwidth]{learning_curves.png}
  \caption{Кривые обучения: loss (слева) и accuracy (справа)}
\end{figure}

\subsection{Анализ}

Начальное значение accuracy (0.627) совпадает с долей класса B в датасете —
это baseline, который даёт модель, предсказывающая всегда класс большинства.
После обучения accuracy на валидации достигла 0.956, что существенно выше baseline.

Наблюдается умеренное переобучение: train loss (0.0421) меньше validation loss
(0.0943) примерно в 2.2 раза, train accuracy (0.9912) превышает validation
accuracy (0.9558) на 3.5 процентных пункта. Кривая валидации выходит на плато
примерно к 30-й эпохе, дальнейшее обучение не даёт прироста.

\section{Прогнозирование}

Команда:

\begin{lstlisting}
python src/predict.py --dataset data_split/data_validation.csv \
                      --model_path saved_model.npy
\end{lstlisting}

\begin{verbatim}
Test samples: 113
Binary cross-entropy: 0.0943
Accuracy: 0.9558

First predictions:
  #00: pred=M true=B [MISS]  p(M)=0.755
  #01: pred=B true=B [ok]    p(M)=0.001
  #02: pred=M true=M [ok]    p(M)=0.994
  #03: pred=M true=M [ok]    p(M)=0.999
  #04: pred=B true=B [ok]    p(M)=0.002
  #05: pred=M true=M [ok]    p(M)=0.999
  #06: pred=B true=B [ok]    p(M)=0.078
  #07: pred=M true=M [ok]    p(M)=1.000
  #08: pred=M true=B [MISS]  p(M)=0.881
  #09: pred=B true=B [ok]    p(M)=0.011
\end{verbatim}

Метрики на валидационной выборке полностью совпали с последней эпохой обучения
(0.0943 и 0.9558), что подтверждает корректность сохранения и загрузки модели,
а также правильное применение нормализации.

Обе ошибочных классификации относятся к типу false positive (модель предсказала
M, тогда как истинный класс B). В контексте медицинской задачи это менее
критично, чем false negative (пропуск злокачественной опухоли), однако всё
равно требует внимания.

Большинство предсказаний уверенные: вероятность $p(M)$ либо меньше 0.1, либо
больше 0.9. Это говорит о хорошем качестве разделяющей границы.

\section{Выводы}

В ходе работы реализован многослойный персептрон с нуля на NumPy. Основные
результаты:

\begin{itemize}
  \item Реализованы прямой и обратный проходы, функции активации, softmax,
        кросс-энтропия, инициализация He, SGD.
  \item Достигнута accuracy 95.6\% на валидационной выборке при 84 эпохах
        обучения.
  \item Модель корректно сохраняется и загружается, включая параметры
        нормализации.
  \item Наблюдается умеренное переобучение, которое можно уменьшить
        регуляризацией (dropout, weight decay) или уменьшением числа
        эпох/нейронов.
  \item Для дальнейшего улучшения можно применить оптимизатор Adam,
        BatchNorm или расширить датасет аугментациями.
\end{itemize}

\end{document}